\documentclass[11pt]{article}
\usepackage[paperwidth=210mm,paperheight=297mm,textwidth=182mm,top=20mm,bottom=22mm]{geometry}
\usepackage{amsmath,amssymb,newtxmath}
\usepackage[no-math]{fontspec}
\setmainfont{Times New Roman}
\usepackage{xeCJK}
\setCJKmainfont[AutoFakeBold=2]{SimSun}
\usepackage{graphicx,booktabs,array,hyperref,placeins}
\hypersetup{hidelinks,unicode=true}
\raggedbottom
\setlength{\emergencystretch}{2em}
\renewcommand{\arraystretch}{1.15}
\usepackage{seqsplit}
\begin{document}
\title{Litz-Wire Modeling and Numerical Verification: Comprehensive Research Report · 03 Round-Wire Analysis and Two-Dimensional Calibration}
\author{Miao Yuxuan\\Huainan Normal University}
\date{}
\maketitle
\noindent 合并稿对应章节 / Chapter in merged manuscript: \href{report_en_single.pdf#chapter-03}{03}\par
\FloatBarrier
\section*{03 Round-Wire Analysis and Two-Dimensional Calibration}
\hypertarget{chapter-03}{}
\pdfbookmark[0]{03 Round-Wire Analysis and Two-Dimensional Calibration}{bookmark-chapter-03}
\subsection*{Round Derivation}
\hypertarget{03.round_derivation.heading}{}
\pdfbookmark[1]{Round Derivation}{bookmark-03.round_derivation.heading}
The axial current in a round wire satisfies (1/r)d(r dJz/dr)/dr=κ²Jz. Regularity at the axis selects I\(_0\)(κr). The identity \(\int\)\(_0\)\(^a\) rI\(_0\)(κr)dr=aI\(_1\)(κa)/κ together with the total-current constraint determines the normalization. The surface axial electric field divided by total current gives the internal impedance; its real part is checked by an independent Joule integral. Thus the derivation connects analytical shape, normalization, and power.


\subsection*{Round-Wire Analysis and Independent Two-Dimensional Calibration}
\hypertarget{03.analytic.heading}{}
\pdfbookmark[1]{Round-Wire Analysis and Independent Two-Dimensional Calibration}{bookmark-03.analytic.heading}
The long straight round-wire internal solution uses modified Bessel functions \href{report_en_single.pdf#E03}{(3)}. At 200 kHz, the skin depth is 0.147771655 mm, producing pronounced surface concentration relative to the 1 mm radius. Frozen analytical values are 5.488101486 mΩ/m DC resistance, 20.016203628 mΩ/m AC resistance, a ratio of 3.647199980, and loss of 8.006481451 W/m. This internal impedance excludes the external inductance determined by the return path. Radial phase retains the original complex phasor and is displayed with a continuous nearest-neighbor 360° lift; unwrapping changes the display, not the underlying data.


Independent two-dimensional disk diffusion uses 353, 1389, and 5573 nodes and is normalized by total complex current. The finest relative AC-resistance error against the analytical reference is 1.79233×10⁻⁶, and the power-integral discrepancy is 3.06800×10⁻⁶. These results support discretization and normalization for the isolated internal-field problem. Its boundary conditions differ from a finite three-dimensional end region, so this error cannot be transferred directly to that problem. The analytical cross-sectional image is formed by linear interpolation of radial analytical magnitudes followed by an axisymmetric mapping.


\hypertarget{E03}{}
\begin{equation}
\begin{gathered}
\delta=(\pi f\mu\sigma)^{-1/2} \\
\kappa=(1+j)/\delta \\
J_z(r)=\frac{I_t\kappa I_0(\kappa r)}{2\pi a I_1(\kappa a)} \\
z_{\mathrm{int}}=\frac{\kappa I_0(\kappa a)}{2\pi a\sigma I_1(\kappa a)}
\end{gathered}
\tag{3}\label{eq:E03}
\end{equation}
\begin{figure}[!htbp]
\centering
\setcounter{figure}{1}
\refstepcounter{figure}\label{fig:q1_profile}
\hypertarget{fig-q1_profile}{}
\includegraphics[width=182mm]{../../../composites/en/q1_profile.pdf}
\par\vspace{4pt}{\small\centering Fig. 2. Round-wire radial current magnitude and phase.\par}
\vspace{3pt}{\small\raggedright The copper conductor has diameter 2 mm and carries 20 A RMS at 200 kHz. (a) RMS magnitude; (b) continuously displayed phase. Phase is displayed with a continuous radial lift, while magnitude retains RMS normalization. Radial samples come from preserved Bessel analytical data.\par}
\end{figure}
\begin{figure}[!htbp]
\centering
\setcounter{figure}{2}
\refstepcounter{figure}\label{fig:q1_cross_section}
\hypertarget{fig-q1_cross_section}{}
\includegraphics[width=88.9mm]{../../../panels/en/q1_cross_section.pdf}
\par\vspace{4pt}{\small\centering Fig. 3. Analytical round-wire cross-sectional magnitude.\par}
\vspace{3pt}{\small\raggedright The copper conductor has diameter 2 mm and carries 20 A RMS at 200 kHz. The color scale represents axial RMS current density. Radial samples are linearly interpolated and mapped axisymmetrically. Coordinates are displayed in millimetres and current density in A/mm².\par}
\end{figure}
\begin{figure}[!htbp]
\centering
\setcounter{figure}{3}
\refstepcounter{figure}\label{fig:q1_fem_error}
\hypertarget{fig-q1_fem_error}{}
\includegraphics[width=88.9mm]{../../../panels/en/q1_fem_error.pdf}
\par\vspace{4pt}{\small\centering Fig. 4. Round-wire two-dimensional calibration errors.\par}
\vspace{3pt}{\small\raggedright The copper conductor has diameter 2 mm and carries 20 A RMS at 200 kHz. Three grids and independent power integrals concern the same internal-field problem. Resistance is referenced to the internal Bessel solution of the same round wire. Power is evaluated by independent integration over the material region.\par}
\end{figure}
\end{document}
